% !TeX root = ../Book3.tex
% 两幅图只画一个坐标邻域；灰色区域没有画出人为外边界。
\begin{tikzpicture}[x=1cm,y=1cm,>=stealth,font=\small]
	\begin{scope}[shift={(-2.5,0)}]
		\fill[blue!9] (0,0) circle (1);
		\draw[thick,blue!70!black] (0,0) circle (1);
		\draw[very thick,blue!70!black,->] (25:1) arc[start angle=25,end angle=88,radius=1];
		\draw[thick,->] (1,0)--(1.7,0) node[right] {\(\nu_B\)};
		\fill (0,0) circle (1.2pt) node[below] {\(p\)};
		\node at (0,.43) {\(B\)};
		\node[below] at (0,-1.25) {\(\partial B\)：球的边界方向};
	\end{scope}
	\begin{scope}[shift={(2.7,0)}]
		\fill[gray!13] (-1.65,-1.15) rectangle (1.65,1.15);
		\fill[white] (0,0) circle (1);
		\draw[thick,red!75!black] (0,0) circle (1);
		\draw[very thick,red!75!black,->] (145:1) arc[start angle=145,end angle=82,radius=1];
		\draw[thick,->] (1,0)--(.37,0) node[below left] {\(\nu_{M_0}\)};
		\fill (0,0) circle (1.2pt) node[below] {\(p\)};
		\node at (0,.43) {挖去的球};
		\node[right] at (1.1,.8) {\(M_0\)};
		\node[below] at (0,-1.25) {\(\partial M_0\)：朝孔洞的外法向};
	\end{scope}
\end{tikzpicture}
